\documentclass[11pt,a4paper]{article}

\usepackage[utf8]{inputenc}
\usepackage[T1]{fontenc}
\usepackage{amsmath,amssymb,amsthm}
\usepackage{hyperref}
\usepackage{listings}
\usepackage{xcolor}
\usepackage{geometry}
\usepackage{booktabs}
\usepackage{enumitem}

\geometry{margin=1in}

\newtheorem{theorem}{Theorem}[section]
\newtheorem{lemma}[theorem]{Lemma}
\newtheorem{definition}[theorem]{Definition}
\newtheorem{remark}[theorem]{Remark}

\lstdefinelanguage{Lean4}{
  morekeywords={def,theorem,lemma,structure,inductive,namespace,end,
    where,open,import,variable,section,match,with,if,then,else,
    deriving,instance,noncomputable,set_option,sorry,by,have,let,
    unfold,cases,simp,omega,obtain,intro,exact,rfl,some,none},
  sensitive=true,
  morecomment=[l]{--},
  morecomment=[n]{/-}{-/},
  morestring=[b]",
  literate={→}{{$\to$}}1 {←}{{$\leftarrow$}}1 {↦}{{$\mapsto$}}1
           {≤}{{$\le$}}1 {≥}{{$\ge$}}1 {≠}{{$\ne$}}1
           {∧}{{$\wedge$}}1 {∨}{{$\vee$}}1 {¬}{{$\neg$}}1
           {∃}{{$\exists$}}1 {∀}{{$\forall$}}1
           {ℕ}{{$\mathbb{N}$}}1 {ℤ}{{$\mathbb{Z}$}}1
           {×}{{$\times$}}1,
}

\lstset{
  language=Lean4,
  basicstyle=\ttfamily\small,
  keywordstyle=\color{blue}\bfseries,
  commentstyle=\color{gray}\itshape,
  stringstyle=\color{red},
  breaklines=true,
  frame=single,
  numbers=left,
  numberstyle=\tiny\color{gray},
  xleftmargin=2em,
  aboveskip=1em,
  belowskip=1em,
}

\title{Formalizing the 3-State Approximate Majority\\Population Protocol in Lean~4}
\author{}
\date{February 2026}

\begin{document}

\maketitle

\begin{abstract}
We present a formalization in Lean~4 (with Mathlib) of the 3-state approximate
majority population protocol introduced by Angluin, Aspnes, and Eisenstat
(2008).  The formalization covers the protocol's state space, deterministic
transition function, configuration representation with a built-in population
size invariant, single-step dynamics, and several key structural invariants
proved by exhaustive case analysis.  We also lay the groundwork for a
probabilistic treatment by defining the uniform random scheduler as a
probability mass function.  All deterministic results are machine-checked
without axioms; the probabilistic components include clearly marked stubs
for future work.
\end{abstract}

\tableofcontents

%% ============================================================
\section{Introduction}
\label{sec:intro}

Population protocols~\cite{AngluinADFP2006} model distributed computation
among a collection of finite-state agents that interact in pairs chosen by
an adversary or random scheduler.  The \emph{approximate majority} protocol
of Angluin, Aspnes, and Eisenstat~\cite{AngluinAE2008} achieves consensus
using only \textbf{three states}---opinion~$x$, blank~$b$, and
opinion~$y$---with expected convergence time $O(n \log n)$ under a uniform
random scheduler.

This document describes a Lean~4 formalization of the protocol's core
definitions, deterministic invariants, and the beginnings of a probabilistic
framework.  The formalization is structured as a Lake project called
\texttt{PopProto} with Mathlib as a dependency.

%% ============================================================
\section{The Protocol}
\label{sec:protocol}

\subsection{States and Transition Rules}

The protocol uses three states:
\[
  Q = \{ x,\; b,\; y \}
\]
where $x$ and $y$ represent opposing opinions and $b$ is a ``blank''
(undecided) state.

When two agents interact, one is designated the \emph{initiator} and the
other the \emph{responder}.  The key asymmetry is:
\begin{quote}
\textbf{One-way property:} The initiator's state never changes.
\end{quote}

The full transition table $\delta(i, r) = (i', r')$ is:

\begin{center}
\begin{tabular}{cc|cc}
\toprule
Initiator $i$ & Responder $r$ & Initiator $i'$ & Responder $r'$ \\
\midrule
$x$ & $x$ & $x$ & $x$ \\
$x$ & $b$ & $x$ & $x$ \\
$x$ & $y$ & $x$ & $b$ \\
$b$ & $x$ & $b$ & $x$ \\
$b$ & $b$ & $b$ & $b$ \\
$b$ & $y$ & $b$ & $y$ \\
$y$ & $x$ & $y$ & $b$ \\
$y$ & $b$ & $y$ & $y$ \\
$y$ & $y$ & $y$ & $y$ \\
\bottomrule
\end{tabular}
\end{center}

Intuitively: a blank responder adopts the initiator's opinion; two agents with
opposing opinions ``neutralize'' the responder to blank; a blank initiator
does nothing.

\subsection{Configurations}

A \emph{configuration} of a population of $n$ agents is a triple
$(x_c, b_c, y_c) \in \mathbb{N}^3$ satisfying $x_c + b_c + y_c = n$.

%% ============================================================
\section{Lean Formalization Strategy}
\label{sec:strategy}

\subsection{Design Decisions}

\begin{table}[h]
\centering
\begin{tabular}{lll}
\toprule
\textbf{Decision} & \textbf{Choice} & \textbf{Rationale} \\
\midrule
Config representation & Struct with proof field & Sum invariant by construction \\
Population size & Type parameter \texttt{Config n} & Type-level separation \\
Counts & \texttt{Nat} (not \texttt{Fin}) & Simpler arithmetic; \texttt{omega} \\
Proof strategy & Exhaustive case split & Only 9 cases; closes by \texttt{rfl}/\texttt{omega} \\
Probability & \texttt{PMF} from Mathlib & Ready-made PMF type \\
\bottomrule
\end{tabular}
\caption{Key design decisions.}
\label{tab:decisions}
\end{table}

\subsection{Project Structure}

The formalization is organized as follows:

\begin{itemize}[nosep]
  \item \texttt{PopProto/State.lean} --- The inductive type \texttt{State}
  \item \texttt{PopProto/Transition.lean} --- The transition function $\delta$
  \item \texttt{PopProto/Config.lean} --- Configuration type with sum constraint
  \item \texttt{PopProto/Step.lean} --- Single-step configuration update
  \item \texttt{PopProto/Invariant/} --- Five invariant proof modules
  \item \texttt{PopProto/Probability/} --- Scheduler, step distribution, Markov chain stub
\end{itemize}

%% ============================================================
\section{Type Definitions}
\label{sec:types}

\subsection{State}

The state type is a simple inductive with three constructors:

\begin{lstlisting}
inductive State
  | x  -- opinion X
  | b  -- blank (undecided)
  | y  -- opinion Y
  deriving DecidableEq, Repr
\end{lstlisting}

We derive \texttt{DecidableEq} (for case splits in \texttt{if} expressions)
and provide a \texttt{Fintype} instance (the universe has exactly 3 elements).

\subsection{Transition Function}

The transition function is defined by pattern matching on all 9 cases:

\begin{lstlisting}
def delta : State $\times$ State $\to$ State $\times$ State
  | (.x, .x) => (.x, .x)
  | (.x, .b) => (.x, .x)
  | (.x, .y) => (.x, .b)
  | (.b, .x) => (.b, .x)
  | (.b, .b) => (.b, .b)
  | (.b, .y) => (.b, .y)
  | (.y, .x) => (.y, .b)
  | (.y, .b) => (.y, .y)
  | (.y, .y) => (.y, .y)
\end{lstlisting}

The one-way property is an immediate consequence:

\begin{lstlisting}
@[simp]
theorem delta_fst (p : State $\times$ State) :
    (delta p).1 = p.1 := by
  rcases p with $\langle$i, r$\rangle$; cases i <;> cases r <;> rfl
\end{lstlisting}

\subsection{Configuration}

A configuration bundles three natural-number counts with a proof that they
sum to the population size~$n$:

\begin{lstlisting}
structure Config (n : $\mathbb{N}$) where
  x_count : $\mathbb{N}$
  b_count : $\mathbb{N}$
  y_count : $\mathbb{N}$
  sum_eq : x_count + b_count + y_count = n
\end{lstlisting}

By carrying \texttt{sum\_eq} as a field, the population-size invariant is
enforced \emph{by construction}: any function that produces a
\texttt{Config n} must supply a proof.  This eliminates an entire class of
potential errors and makes the invariant trivially true (Theorem~\ref{thm:popsize}).

Helper definitions include:
\begin{itemize}[nosep]
  \item \texttt{gap : $\mathbb{Z}$} --- the signed difference $x_c - y_c$
  \item \texttt{opinionated : $\mathbb{N}$} --- the count $x_c + y_c$
  \item \texttt{allX}, \texttt{allY}, \texttt{allB} --- consensus predicates
  \item \texttt{countOf : State $\to$ $\mathbb{N}$} --- state-indexed lookup
\end{itemize}

%% ============================================================
\section{Step Function}
\label{sec:step}

The function \texttt{Config.step} applies one interaction to a configuration.
Given an initiator state~$i$ and responder state~$r$, it:
\begin{enumerate}[nosep]
  \item Checks feasibility (enough agents of the required types).
  \item Computes the new responder state $r' = \delta(i,r)_2$.
  \item Adjusts the counts: $\mathrm{count}(r)$ decreases by 1 and
        $\mathrm{count}(r')$ increases by 1 (when $r \neq r'$).
  \item Returns the new configuration wrapped in \texttt{Option}, or
        \texttt{none} if infeasible.
\end{enumerate}

The function is defined by explicit pattern matching on all 9 cases.
For each non-trivial case, the proof obligation
$x' + b' + y' = n$ is closed by \texttt{omega} using the hypothesis
\texttt{c.sum\_eq} and the feasibility guard.

A total wrapper \texttt{stepOrSelf} returns the original configuration
when the interaction is infeasible:
\begin{lstlisting}
def stepOrSelf (c : Config n) (i r : State) : Config n :=
  (c.step i r).getD c
\end{lstlisting}

%% ============================================================
\section{Invariant Proofs}
\label{sec:invariants}

All invariant proofs follow the same pattern:
\begin{enumerate}[nosep]
  \item Unfold \texttt{step} to expose the 9-way match.
  \item Apply \texttt{cases i <;> cases r} to split into 9 goals.
  \item Close trivial goals with \texttt{simp\_all} and \texttt{omega}.
  \item For non-trivial goals, destructure the \texttt{Option.some}
        hypothesis and simplify.
\end{enumerate}

\subsection{Population Size}
\label{sec:popsize}

\begin{theorem}[Population Size Invariant]
\label{thm:popsize}
For any configuration $c$ of size $n$ and any interaction $(i,r)$, if
$c' = \texttt{step}(c, i, r)$, then $c'.x + c'.b + c'.y = n$.
\end{theorem}

This is trivially true because \texttt{Config n} enforces
$x + b + y = n$ by construction.  The proof is simply \texttt{c'.sum\_eq}.

\subsection{One-Way Property}

\begin{theorem}[One-Way Property]
For any states $i, r$, we have $\delta(i,r)_1 = i$.
\end{theorem}

Proved by \texttt{cases i <;> cases r <;> rfl} --- each of the 9 cases
is definitionally equal.

\subsection{Gap Bounded Change}

\begin{theorem}[Gap Invariant]
\label{thm:gap}
If $c \xrightarrow{(i,r)} c'$, then $|u' - u| \leq 1$ where
$u = x_c - y_c$ and $u' = x_{c'} - y_{c'}$ (computed over $\mathbb{Z}$).
\end{theorem}

This is proved by exhaustive 9-case analysis. In each case, at most one
of $x_c$ or $y_c$ changes by~$\pm 1$, so the gap changes by at most~1.

\subsection{Absorbing States}

\begin{theorem}[Absorbing Configurations]
The configurations all-$x$, all-$y$, and all-$b$ are absorbing:
if $c$ is one of these, then $\texttt{stepOrSelf}(c, i, r) = c$ for all
$i, r$.
\end{theorem}

When all agents share the same state, every interaction is between two
agents of that state, and $\delta(s,s) = (s,s)$ for all $s$.

\begin{theorem}[Unreachability of All-Blank]
\label{thm:allb}
If $c$ has at least one opinionated agent ($x_c + y_c > 0$), then for any
step $c \to c'$, the configuration $c'$ also has at least one opinionated
agent.  In particular, all-$b$ is unreachable.
\end{theorem}

The proof uses the \texttt{hasOpinion\_preserved} lemma: the last non-blank
agent can only interact with blanks, and in every such interaction (e.g.,
$(x, b) \to (x, x)$), the blank is converted rather than the opinionated
agent being eliminated.

\subsection{Last Agent Survival}

\begin{theorem}[Last Agent]
If $x_c \geq 1$ and $y_c = 0$, then after any step, $x_{c'} \geq 1$ and
$y_{c'} = 0$.  Symmetrically for~$y$.
\end{theorem}

With no opposing agents, the only possible interactions involving an $x$-agent
are $(x,x)$ and $(x,b)$, both of which preserve or increase the $x$-count.

%% ============================================================
\section{Probabilistic Framework}
\label{sec:probability}

\subsection{Uniform Random Scheduler}

Under the uniform random scheduler, an ordered pair of distinct agents
$(a_i, a_j)$ is selected uniformly from $n(n-1)$ possibilities.
The probability of an interaction of type $(s_1, s_2)$ is:
\[
  \Pr[(s_1, s_2)] =
  \begin{cases}
    \dfrac{c(s_1) \cdot (c(s_1) - 1)}{n(n-1)} & \text{if } s_1 = s_2, \\[6pt]
    \dfrac{c(s_1) \cdot c(s_2)}{n(n-1)} & \text{if } s_1 \neq s_2.
  \end{cases}
\]

We define \texttt{interactionCount} and \texttt{interactionProb} in
\texttt{Scheduler.lean}.  The key identity
\[
  \sum_{s_1 \in Q} \sum_{s_2 \in Q}
  \texttt{interactionCount}(c, s_1, s_2) = n(n-1)
\]
is stated as \texttt{sum\_interactionCount}.  This identity follows from
$(x+b+y)^2 - (x+b+y) = n^2 - n = n(n-1)$, but its formal proof requires
careful handling of natural-number subtraction (since $a \cdot (a-1)$ in
$\mathbb{N}$ truncates to~0 when $a=0$).  This proof is currently left as
a \texttt{sorry} stub.

The PMF construction \texttt{interactionPMF} wraps the probabilities into
Mathlib's \texttt{PMF (State $\times$ State)} type.

\subsection{One-Step Distribution}

The one-step distribution \texttt{stepDist} composes the scheduler PMF
with the deterministic step function via \texttt{PMF.map}:

\begin{lstlisting}
noncomputable def stepDist (c : Config n)
    (hn : n $\ge$ 2) : PMF (Config n) :=
  PMF.map (fun p => c.stepOrSelf p.1 p.2)
          (c.interactionPMF hn)
\end{lstlisting}

\subsection{Markov Chain (Stub)}

The file \texttt{MarkovChain.lean} outlines how to view the protocol as a
discrete-time Markov chain using Mathlib's
\texttt{ProbabilityTheory.Kernel} framework.  Constructing the transition
kernel and proving convergence properties are left as future work.

%% ============================================================
\section{Proof Techniques}
\label{sec:techniques}

The formalization relies heavily on a small set of Lean~4 / Mathlib tactics:

\begin{description}[nosep]
  \item[\texttt{cases i <;> cases r}]
    Splits into 9 subgoals, one per $(i,r)$ pair.  Most transition-level
    properties reduce to definitional equality after this split.

  \item[\texttt{simp\_all}]
    Simplifies goals using hypotheses and \texttt{@[simp]} lemmas.
    Particularly effective after case-splitting, when the match arms
    of \texttt{step} have been reduced.

  \item[\texttt{omega}]
    Decides linear natural-number and integer arithmetic.  Closes
    sum-preservation goals, count bounds, and gap-change bounds.

  \item[\texttt{split\_ifs}]
    Case-splits on \texttt{if}-\texttt{then}-\texttt{else} expressions
    remaining after partial simplification.

  \item[\texttt{obtain}]
    Destructures existential hypotheses produced when \texttt{simp\_all}
    processes the \texttt{Option.some} constructor in step results.
\end{description}

%% ============================================================
\section{Challenges and Future Work}
\label{sec:challenges}

\subsection{Challenges Encountered}

\begin{enumerate}
  \item \textbf{Natural-number subtraction.}  The expression $a \cdot (a-1)$
        in $\mathbb{N}$ introduces truncation artifacts.  The tactic
        \texttt{omega} handles linear arithmetic over $\mathbb{N}$ well,
        but the nonlinear identity $\sum c_i(c_i - 1) + \sum_{i \neq j}
        c_i c_j = n(n-1)$ required for the scheduler's sum-to-1 proof
        resists both \texttt{omega} (nonlinear) and \texttt{nlinarith}
        (Nat subtraction).  Converting via \texttt{zify} is a promising
        approach but requires careful handling of \texttt{Nat.cast} applied
        to subtraction terms.

  \item \textbf{Struct field unfolding.}  When \texttt{step} returns
        \texttt{some \{x\_count := \ldots, \ldots\}}, the hypothesis after
        \texttt{simp\_all} takes the form $\exists h, \{\ldots\} = c'$,
        requiring \texttt{obtain} to extract the proof and substitute.

  \item \textbf{Linter noise.}  Automated proofs using chains like
        \texttt{cases i <;> cases r <;> simp\_all <;> \ldots} produce
        linter warnings about unused or unreachable tactics on branches
        that close early.  We suppress these with targeted
        \texttt{set\_option} directives.
\end{enumerate}

\subsection{Future Work}

\begin{enumerate}
  \item \textbf{Complete the scheduler sum-to-1 proof.}  The algebraic
        identity is straightforward on paper; formalizing it requires
        either working in $\mathbb{Z}$ (via \texttt{zify}) or establishing
        auxiliary lemmas about $a \cdot (a-1)$ in $\mathbb{N}$.

  \item \textbf{Construct the Markov transition kernel.}  Lift
        \texttt{stepDist} into Mathlib's \texttt{ProbabilityTheory.Kernel}
        framework by equipping \texttt{Config n} with the discrete
        $\sigma$-algebra.

  \item \textbf{Convergence to consensus.}  Prove that, starting from any
        configuration with $x_c + y_c > 0$, the protocol reaches all-$x$
        or all-$y$ with probability~1.

  \item \textbf{Expected convergence time.}  Formalize the $O(n \log n)$
        bound from~\cite{AngluinAE2008} using a potential-function argument
        and optional stopping times.

  \item \textbf{Correctness of majority.}  Show that if the initial
        majority is $x$ (i.e., $x_0 > y_0$), then the protocol converges
        to all-$x$ with probability $1 - o(1)$ as $n \to \infty$.
\end{enumerate}

%% ============================================================
\begin{thebibliography}{9}

\bibitem{AngluinADFP2006}
D.~Angluin, J.~Aspnes, Z.~Diamadi, M.~J.~Fischer, and R.~Peralta.
\newblock Computation in networks of passively mobile finite-state sensors.
\newblock \emph{Distributed Computing}, 18(4):235--253, 2006.

\bibitem{AngluinAE2008}
D.~Angluin, J.~Aspnes, and D.~Eisenstat.
\newblock A simple population protocol for fast robust approximate majority.
\newblock \emph{Distributed Computing}, 21(2):87--102, 2008.

\end{thebibliography}

\end{document}
